\subsection{Transaction cost models as inputs to portfolio construction}

\begin{qfmodel}[Transaction cost functions for portfolio construction: flat, linear, piecewise-linear and quadratic forms, with units and size conventions]\label{mod:portfolio-transaction-cost-function}
\qffield{Statement} Narang's transaction cost model quantifies ``the cost of making a trade of a given size'' so that the portfolio construction model can weigh it against the alpha and risk models when choosing the portfolio to hold \citep[chap.~5, ``Transaction Cost Models'', p.~85]{s7-narang}. It is ``not designed to minimize the cost of trading, only to inform the portfolio construction engine''; minimising cost is the execution algorithm's job (p.~86), and cost is reduced in two phases, first by the portfolio construction model using the cost estimate and then by execution (sec.~5.3, p.~96). In his framing there are ``only two reasons to make a trade'', a better reward from the alpha model or a lower risk from the risk model, and a trade's benefit must ``clear the hurdle of the cost of transacting'' (p.~85).

He names four forms, each giving total cost against trade size, with commissions and fees as the baseline (sec.~5.2, pp.~91--96, Exhibits~5.1--5.4):
\begin{itemize}
\item \emph{flat}: the same cost whatever the size; reasonable when the size traded is nearly always about the same and liquidity is fairly constant, ``rarely correct'' otherwise (pp.~92--93);
\item \emph{linear}: cost rising with a constant slope in size, overestimating at small sizes and underestimating at large ones in his exhibit (pp.~93--94);
\item \emph{piecewise-linear}: linear in ranges, with a steeper line beyond each breakpoint; ``rather popular among quants as a happy medium between simplicity and accuracy'' (p.~94);
\item \emph{quadratic}: cost rising ever faster with size; in his words the most ``accurate'' of the four in the exhibits and the most computationally intensive (p.~95).
\end{itemize}

\qffield{Notation} Interpretation, for a single equity traded in shares: let $q$ be the signed trade size in shares ($q>0$ a buy, $q<0$ a sale), $V>0$ a volume scale in shares such as average daily volume, $P>0$ the price in currency per share, and $C(q)$ the \emph{total} expected cost of the trade in currency. For $q\neq0$ the \emph{per-share} cost is $c(q)=C(q)/\lvert q\rvert$ in currency per share, and $c(q)/P$ is the same cost as a fraction of price (times $10^4$ for basis points); $c(0)$ is not defined. With coefficients $a\ge0$ in currency, $b,b_0,h_i\ge0$ in currency per share, $0<\kappa_1<\dots<\kappa_m$ in shares, $k\ge0$ in currency per share squared and $\eta\ge0$ dimensionless, the forms are
\[ C_{\mathrm{flat}}(q)=\begin{cases}a,&q\neq0,\\0,&q=0,\end{cases}\qquad C_{\mathrm{lin}}(q)=b\lvert q\rvert,\qquad C_{\mathrm{quad}}(q)=b\lvert q\rvert+kq^2, \]
\[ C_{\mathrm{pwl}}(q)=b_0\lvert q\rvert+\sum_{i=1}^{m}h_i\bigl(\lvert q\rvert-\kappa_i\bigr)_{+},\qquad C_{\mathrm{pow}}(q)=\eta P\lvert q\rvert\Bigl(\frac{\lvert q\rvert}{V}\Bigr)^{1/2}, \]
with $(x)_+=\max(x,0)$. Every form has $C(0)=0$ and $C\ge0$, and depends on $q$ only through $\lvert q\rvert$. All but the flat form are continuous; the piecewise-linear form has slope $b_0+\sum_{i:\kappa_i<\lvert q\rvert}h_i$, rising at each breakpoint $\kappa_i$, and equals $\max_j(s_j\lvert q\rvert-d_j)$ with $s_j=b_0+\sum_{i\le j}h_i$ and $d_j=\sum_{i\le j}h_i\kappa_i$ ($s_0=b_0$, $d_0=0$). A maximum of lines with increasing slopes alone does not give $C(0)=0$: its intercepts must take this form.

For $b_0=0.01$, $h=(0.02,0.05)$ and $\kappa=(10{,}000,50{,}000)$: $C=0$ at $q=0$; 50, 100, 700, 1{,}300 and 2{,}100 at $\lvert q\rvert=5{,}000$, $10{,}000$, $30{,}000$, $50{,}000$ and $60{,}000$, the same for buys and sales; continuous at both breakpoints; equal to the maximum of lines with slopes 0.01, 0.03, 0.08 and intercepts 0, 200, 2{,}700.

The forms assume a cost symmetric between buys and sales, never negative (no rebates or favourable price drift), expected rather than realised, for one instrument and one period, independent of the trade's duration and of other trades. Contracts or currency notional need their own size and price units, such as a contract multiplier; the per-share conversions above do not carry over unchanged.
A total cost linear in size is a constant per-share cost; a quadratic total cost is a per-share cost linear in size, $c(q)=b+k\lvert q\rvert$ for $q\neq0$. A per-share commission of, say, \$0.001 (Narang's example, p.~87) is therefore a linear, not a flat, cost; only a charge per ticket is flat. A per-share cost proportional to $(\lvert q\rvert/V)^{1/2}$, as in \hyperref[ev:square-root-impact]{the square-root law}, is a total cost proportional to $\lvert q\rvert^{3/2}$.

Numerical illustration: with $C_{\mathrm{pow}}$, $\eta=0.5\times2\%=0.01$, $P=\$50$ and $V=10^6$ shares, 10{,}000 shares bought or sold cost 10~bp per share, \$500 in total; 20{,}000 shares cost 14.1~bp, \$1{,}414; doubling the size multiplies the per-share cost by $2^{1/2}$ and the total by $2^{3/2}\approx2.83$. Using the peak-impact law as a cost scale is a simplification: \hyperref[ev:square-root-impact]{the evidence entry} gives an execution shortfall per share of about two-thirds of peak impact.

\qffield{Conditions} The ranking in the schematic comparison is not an established result. Narang's exhibits are schematic, with unlabelled axes and no data; the ``true'' cost curve in each is drawn convex and starting at zero (Exhibits~5.1--5.4, pp.~93--95), although his text puts fixed charges under it as a floor (p.~91). Against an assumed convex curve, the closeness of each form follows from the shape drawn, not from estimation. His own summary is conditional: simpler models suit traders whose sizes are roughly constant or small, more complex ones those who may trade significant or variable quantities quickly, and ``any of the four models described here can be valid in the right set of circumstances'' (sec.~5.3, p.~96).

His statement that a quadratic shape ``is generally accepted by the quant community'' (p.~92) is unsourced, and his note~2 qualifies it: ``there is good evidence that, at least for U.S. equities, cost per share scales with an exponent of 1.5, or in other words, as the square root of trade size'' (n.~2, p.~97). That sentence conflates two conventions: a per-share cost scaling as the square root of size is a total cost with exponent 1.5 (Notation).

Both a quadratic and a 3/2-power total cost are convex, so ``the cost gets higher ever more quickly'' as size grows (p.~92) does not distinguish them; which exponent fits is an empirical question (\hyperref[ev:square-root-impact]{square-root impact}, \hyperref[meth:almgren-impact-estimation]{direct estimation}, whose fitted temporary-impact exponent of 0.6 per share implies, at a fixed duration, a total exponent of about 1.6).

\qffield{Practice} Narang describes models built separately for each instrument, ``highly empirical'', evolving with the quant's own recorded executions and refreshed to keep up with market conditions; a less liquid small-cap stock costs more to trade than a mega-cap (pp.~91--92 and~96). His Exhibit~5.4 contrasts the expected cost, specified before trading, with what is ``observed empirically after trading''; the differences come from changes in liquidity, volatility and the mix of traders (p.~95). Interpretation: the realised cost observed after trading is measured against a benchmark price, and what is observed is a cost, not the impact component; how ex-ante cost forecasts are tested against realised costs is in \hyperref[meth:transaction-cost-analysis]{transaction cost analysis}.

Interpretation: inside an optimiser the cost enters the objective jointly with the alpha forecast and the risk penalty, so a trade is not tested against a separate hurdle one by one; and trades can also be forced by flows, margin or constraints, which Narang's two reasons leave out. The cost of underestimating or overestimating costs (too much turnover, or positions held too long, p.~86) is the subject of \hyperref[pit:backtest-cost-short-assumptions]{backtest cost and short assumptions}.

\qffield{Pitfalls} \emph{No universal ranking of accuracy or computation.} Narang's comparisons of accuracy and computation rest on his drawn curve and do not hold generally. Which form is most accurate depends on the instrument, the size range traded and the true shape, which Narang's own note~2 says may not be quadratic. Computational cost depends on how the form enters the optimiser, not on how many terms it has.

What the solver needs: Boyd et al.'s generic separable cost per asset is $x\mapsto a\abs{x}+b\sigma\abs{x}^{3/2}/V^{1/2}+cx$ for a dollar trade $x$, with $a$ half the spread plus linear fees, $b$ a constant fitted to realised costs (``around one'' by the rule of thumb that a day's volume moves the price by a day's volatility), $V$ the period's dollar volume and $c$ an asymmetry term, which makes the cost negative when $\abs{c}>\abs{a}$; they know of no empirical test of this exact form, though similar models with cost growing about as the 3/2 power of size are validated in the work they cite \citep[sec.~2.3, eq.~(2.2), pp.~9--10]{exp-boyd-multiperiod-trading}. Piecewise-linear and added quadratic terms are common variants and ``almost all of these are convex functions''; a fixed fee for any nonzero trade is not (p.~11).

The optimisation is fast and reliable so long as costs, risk and constraint sets are convex (the self-financing equality relaxed to an inequality, tight at the optimum); a single-period problem with a factor risk model costs $O(nk^2)$ flops with a leading coefficient of order 100 by interior-point methods, ``even for \dots\ the 3/2 power transaction costs'', and a purely quadratic objective with only linear equality constraints has a coefficient closer to 2, about 50 times faster (sec.~4.7, pp.~38--39). A multi-period problem over $H$ periods has $Hn$ variables, but its structure lets the effort grow linearly in $H$ (sec.~5.3, p.~51). So the computational difference that matters is convex against nonconvex (flat per-trade fees, integer and minimum-size constraints) and quadratic against general convex, not the number of terms.

Interpretation, not stated by Boyd et al.: a convex piecewise-linear cost can be written with linear constraints and auxiliary variables. Narang's own wording is hedged (``this perhaps more accurate one is more computationally intensive'', n.~2, p.~97). For example, Boyd et al.\ use a spread term plus a $\lvert z\rvert^{3/2}$ impact term inside a convex problem \citep[author version of 5 Jan.\ 2024, sec.~2.4, eq.~(10), p.~17]{wide-boyd-markowitz-seventy}, and find integer, cardinality and minimum-size constraints nonconvex and much slower to solve (sec.~4.3, pp.~27--28); see \hyperref[mod:mean-variance-portfolio-optimisation]{mean--variance optimisation}.

\emph{Units.} A cost function stated without its convention (total or per unit; shares, currency or fraction of volume) cannot be compared with an impact law; say which. \emph{Benefit and cost.} The trade-off is between expected benefit and expected cost over the same horizon; comparing a forecast return with a cost in different units or horizons misstates the hurdle.

\emph{Single- against multi-period optimisation}. Boyd et al.\ plan trades over a horizon $H$ and execute only the first period's. They argue that this helps when forecasts decay, when risk, constraints or liquidity are known to change, and when a portfolio is being set up or shut down: ``trading towards a lower risk position starts well before the increase in risk, trading it off with the transaction costs''. Their example with $H=2$ shows a ``substantial advantage'' over single-period optimisation, but they attribute it ``mostly'' to the extra forecast of tomorrow's return, and their return forecasts are built from realised returns plus noise, which they call ``clearly not a practical forecast'' \citep[secs.~5.2 and~7.4, pp.~46--47 and~66--69]{exp-boyd-multiperiod-trading}. The example therefore shows the mechanics, not a measured gain.

\qffield{See also} \hyperref[def:transaction-costs]{transaction costs} (the components, including Narang's commissions and fees, slippage and impact); \hyperref[ev:square-root-impact]{square-root impact} and \hyperref[meth:almgren-impact-estimation]{direct estimation of impact} (the empirical shape); \hyperref[mod:almgren-chriss-execution]{mean--variance execution} (linear temporary impact, a quadratic total cost for a given schedule); \hyperref[meth:transaction-cost-analysis]{transaction cost analysis} (forecasts against realised costs); \hyperref[pit:backtest-cost-short-assumptions]{backtest cost and short assumptions}.

\qffield{Further reading} Guasoni and Muhle-Karbe's survey summarises the literature \citep[arXiv v1, sec.~1, PDF p.~3]{guasoni-muhle-karbe-2012}. With proportional costs, ``optimal portfolios entail a no-trade region, in which it is optimal to keep existing holdings'', trading happens only at its boundaries, and the region ``is wide, even for small transaction costs'' while its welfare cost is small (Constantinides, 1986). Davis and Norman solve the free-boundary problem for power utility from consumption over an infinite horizon. The width of the no-trade interval and the utility loss are of order $\varepsilon^{1/3}$ and $\varepsilon^{2/3}$ in the proportional cost $\varepsilon$ (Shreve and Soner; Rogers). The survey derives the leading constant of the width for its long-run model (Lemma~3.2 and the expansion that follows it, PDF p.~9).

\emph{Davis and Norman through Shreve and Soner}. Shreve and Soner treat the same problem, one money market and one stock, and describe Davis and Norman's paper as providing ``a precise formulation and analysis, including an algorithm and numerical computations of the optimal policy'', but ``under restrictive and not fully verifiable assumptions'', which they replace ``by the sole assumption of a finite value function'', confirming some of Davis and Norman's conjectures and refuting others. Their result: ``regardless of model parameters, the proportion of wealth held in the stock by the optimal portfolio remains in an interval whose endpoints depend on all the model parameters, but not on the wealth''; when the Merton proportion exceeds one (leverage is optimal) it can lie outside the interval, so ``the presence of transaction costs reduces the agent's desire for leverage''.

They also show in their appendix that the liquidity premium behaves like the transaction cost to the $2/3$ power in this model \citep[\emph{Ann.\ Appl.\ Probab.} 4(3), 1994, sec.~1, pp.~610--612]{shreve-soner-1994}. That supports the wedge description above with the proof's conditions stated, from a source that read Davis and Norman, not from their paper. These are restated results, not Davis and Norman's own theorem. Interpretation: the rebalancing rule this implies, trading to the nearest boundary only when the position leaves the region, differs from penalising turnover in a single-period optimiser, which is the form in Boyd et al.

\end{qfmodel}
